\section{General bundle formulas and additive invariants}
\label{sec:general-operations}
The results in this appendix use the general projective-bundle theorem of
Iritani--Koto. For product computations, the cubic proof and the classification and Hodge
theorems use only Lemma~\ref{lem:special-projective-line} and the ruled-product
calculation.

\paragraph{Coefficient fields.}
The intrinsic ample-adic completion and the Laurent neighborhood of a
fibre or exceptional cusp need not map to one another.  Iritani's
display~(1.1) supplies the required common coefficient maps
\cite[(1.1)]{IritaniBlowup}; injectivity for each center summand is
proved in Lemma~\ref{lem:faithful-center-base-change}.  For projective bundles the common source is
\(\C[z,q][[Q,\widehat\tau]]\).  On the blowup side the ambient term uses the
corresponding graded source.  A center term begins with its intrinsic reduced
coefficient ring from Remark~2.3, maps to the possibly smaller image ring of
Remark~5.6, and is then pulled to a separate target ring for that summand in
the external direct sum.  Embed the generic fields of the ambient common image
and of these separate center images in a fixed algebraically closed overfield
and form every primary block there.  Base
change after this splitting extension preserves the rank, nilpotent Jordan
type, regular formal gauge, and formal exponent classes of each block; forming
primary factors over a nonsplitting field before extension would not.

\paragraph{Restriction to even cohomology.}
The comparison is restricted to even cohomology, not merely to the
even bulk base.  The supermodule conventions and the pull-push and Fourier
formulas in Iritani's Section~5.8 and Iritani--Koto's construction are
parity-preserving.  After the odd bulk variables are set to zero, the
comparison and its inverse therefore restrict to mutually inverse maps on
\(H^{\mathrm{ev}}\).  This is the QDM on which the additive invariant is defined.

\paragraph{Regularity in \(z\).}
The comparisons are regular in \(z\).  Iritani's Theorem~5.18 and
Iritani--Koto's Theorem~5.1 are respectively defined over rings of the form
\[
 \C[z]((q_{\mathrm{exc}}^{-1/\mathfrak s}))[[Q,\widetilde\tau]],
 \qquad
 \C[z]((q_{\mathrm{fib}}^{-1/r'}))[[Q,\widehat\tau]],
\]
where Iritani--Koto take \(r'=r\) when \(r-1\) is even and \(r'=2r\) when
\(r-1\) is odd.  Here \(\mathfrak s\in\{c-1,2(c-1)\}\) is Iritani's
parity-dependent ramification index from \cite[(5.11)]{IritaniBlowup}.  Their
inverses use no negative powers of \(z\); see also
\cite[Remark~5.3]{IritaniKoto}.  The maps are homogeneous but need not have
degree zero: the projective-bundle map has degree \(-(r-1)\), while Iritani's
ambient and center components have degrees \(0\) and \(-c\), respectively
\cite[Theorem~5.1(3)]{IritaniKoto}\cite[Theorem~5.18(3)]{IritaniBlowup}.
The cited maps intertwine the actual \(z\)-connections.  A common grading
normalization can translate the residue by a scalar, preserving its
discriminant; no eigenline is independently regraded.  The graded completed rings embed coefficientwise
in \(\overline K[[z]]\); no identification with an ordinary tensor-product
completion is needed.

\begin{proposition}[Intrinsic formulas for the two counts]
\label{prop:intrinsic-count-formulas}
\coverage{absent}%
\uses{lem:faithful-center-base-change,prop:rank2-rigidity}%
\imports{IritaniKoto:thm-5.1,IritaniBlowup:thm-5.18}%
For either \(I=I_{\mathrm{exp}}\) or \(I=I_{\mathrm{lat}}\), a smooth
codimension-\(c\) center \(Z\subset Y\) and a rank-\(r\) vector bundle \(V\)
satisfy
\begin{equation}
 I(\Bl_ZY)=I(Y)+(c-1)I(Z),\qquad
 I(\PP_Y(V))=rI(Y).
 \label{eq:intrinsic-count-formulas}
\end{equation}
The formulas are componentwise for disconnected centers.
\end{proposition}

\begin{proof}
\emph{Projective bundles.}
For a rank-\(r\) bundle, Iritani--Koto's Theorem~5.1 gives \(r\) copies of
the base QDM after Laurent extension \cite[Theorem~5.1]{IritaniKoto}.  Its
invertible joint bulk Jacobian gives independent unit coordinates on the
\(r\) target factors.  In the \(j\)-th factor, changing its unit coordinate
adds an independent scalar \(u_jI\) to Euler multiplication.  For two finite
spectra, the resultant of their translated characteristic polynomials is a
nonzero polynomial in \(u_i-u_j\), with leading coefficient one.  Hence the
spectra of distinct factors are disjoint at the generic point: their primary
blocks do not fuse when the direct sum is formed.  Regularity, restriction to
even cohomology, the fixed splitting field, and invariance under common scalar grading shifts therefore
carry the invariant blockwise across the comparison, and additivity gives
\eqref{eq:intrinsic-count-formulas}.  Tensoring the bundle by a line bundle, when needed
for the global-generation convention, changes neither the projective bundle
nor the intrinsic Novikov embedding
\cite[Remark~1.2 and Remark~5.2]{IritaniKoto}.

\smallskip
\noindent\emph{Blowups.}
For a codimension-\(c\) blowup, Iritani's Theorem~5.18 gives one ambient QDM
and \(c-1\) center QDMs \cite[Theorem~5.18]{IritaniBlowup}.  Parts (4), (6),
and (7) give the comparison asymptotics, the leading reconstruction
parameters, and the invertible combined bulk Jacobian, respectively.  The
independent unit coordinates supplied by part~(7) give the same nonzero
resultant argument, so the spectra of the ambient and individual center
summands are pairwise disjoint generically.  Regularity, restriction to even
cohomology, the fixed splitting field, invariance under common scalar grading shifts, and
Lemma~\ref{lem:faithful-center-base-change} for each center copy then give
\eqref{eq:intrinsic-count-formulas}.  Faithful center base change identifies each of the \(c-1\) contributions
with \(I(Z)\).  Each comparison is evaluated over its own common
coefficient ring; the proof never composes Laurent comparisons recursively.
\uses{lem:faithful-center-base-change,prop:rank2-rigidity}%
\imports{IritaniKoto:thm-5.1,IritaniBlowup:thm-5.18}%
\end{proof}

\subsection{The additive formulation}
For retained block data \(\mathcal T\), let \(\mathcal L_{\mathcal T}\)
be the monoid of finite multisets of regular-isomorphism classes of blocks,
with addition by multiset union.  Write \(\mathcal B_{\mathcal T}(Y)\) for the
block multiset of \(Y\).  An \emph{additive invariant} is a homomorphism
\[
 w:\mathcal L_{\mathcal T}\longrightarrow A,
 \qquad I_w(Y):=w(\mathcal B_{\mathcal T}(Y)),
\]
where \(A\) is a commutative monoid, possibly without subtraction.  For a blowup
\(\beta:\Bl_ZY\to Y\) of codimension \(c\ge2\), retain the data of each
center occurrence separately:
\[
 \omega_j=(\beta,Z,j,\varsigma_j,\chi_j),
 \qquad 0\le j\le c-2.
\]
Here \(\varsigma_j\) is the reconstruction parameter and \(\chi_j\) its
numerical Novikov specialization.  Set
\(c_w(\omega_j):=w(\mathcal B_{\mathcal T}(Z;\omega_j))\);
equality of these values need not identify the block multisets.

\begin{theorem}[A birational invariance criterion]
\label{thm:marker-ledger}
\coverage{conditional_deduction}%
\lean{effectiveBlockLedger_fold_unique,%
  occurrenceIndexedMarker_eq_of_birational,%
  occurrenceIndexedMarker_descends_to_birationalClasses}%
\imports{AKMW:weak-factorization}%
Let \(w:\mathcal L_{\mathcal T}\to A\) be an additive invariant unchanged by the
allowed regular scalar extensions and formal coordinate changes.  Suppose
\begin{equation}
 I_w(\Bl_ZY)=I_w(Y)+\sum_{j=0}^{c-2}c_w(\omega_j),
 \qquad\operatorname{codim}_YZ=c\ge2.\label{eq:marker-blowup}
\end{equation}
If every summand associated with a smooth center of dimension at most \(d-2\)
has zero correction, then \(I_w\) is a birational invariant of smooth
projective \(d\)-folds.
\end{theorem}

\begin{proof}
Projective weak factorization writes a birational map of smooth projective
\(d\)-folds as a zigzag of blowups and blowdowns along smooth centers of
codimension at least two \cite[Theorem~0.3.1]{AKMW}.  By
\eqref{eq:marker-blowup} and the vanishing of every center correction, the
invariant is unchanged at each step; the equalities therefore telescope.
Disconnected centers are treated componentwise, while a divisor blowup is an
isomorphism.
\end{proof}

For the quantum comparisons one also has the bundle formula
\begin{equation}
 I_w(\PP_Y(V))=rI_w(Y),\qquad\rk V=r.
 \label{eq:marker-pbundle}
\end{equation}
It is not a hypothesis of the birational invariance criterion.

\begin{proposition}[Marker formulas for the QDM comparisons]
\label{prop:qdm-operation-ledgers}
\coverage{conditional_deduction}%
\lean{qdmComparison_evenCarrierEquiv,%
  qdmProjectiveBundle_markerFormula_of_suspensionComparison,%
  qdmBlowup_markerFormula_of_suspensionComparison,%
  qdmProjectiveBundle_markerFormula_of_blockComparison,%
  qdmBlowup_markerFormula_of_occurrenceBlockComparison,%
  qdmBlowup_step_of_occurrenceBlockComparison}%
\imports{IritaniKoto:thm-5.1,IritaniBlowup:thm-5.18}%
Let an additive invariant be preserved by the regular scalar extensions and formal
coordinate changes described above, and by every recorded grading suspension
if its block type retains grading.  The projective-bundle and blowup QDM
comparisons on even cohomology then supply \eqref{eq:marker-pbundle} and
\eqref{eq:marker-blowup}; the contribution of the \(j\)-th center summand is
evaluated using its separate data \(\omega_j\).
\end{proposition}

\begin{proof}
The comparison proof of Proposition~\ref{prop:intrinsic-count-formulas}
uses only regular transport, generic separation and additivity.
Apply it to \(w\), retaining each center occurrence and any grading
suspensions in its block data.
\uses{prop:intrinsic-count-formulas}%
\imports{IritaniKoto:thm-5.1,IritaniBlowup:thm-5.18}%
\end{proof}

\begin{corollary}[Threefold criterion]
\label{cor:threefold-marker}
\coverage{fragment}%
\lean{threefold_not_rational_of_occurrenceIndexedMarker,%
threefold_projectiveLine_irrationality_of_positive_count}%
Let \(Y\) be a smooth projective complex threefold.  If
\(I_{\mathrm{lat}}(Y)>0\), then \(Y\times\PP^1\) is irrational, and so
is \(Y\).  In particular, this holds when \(I_{\mathrm{exp}}(Y)>0\).
\end{corollary}

\begin{proof}[Proof of Corollary~\ref{cor:threefold-marker}]
The projective-bundle formula gives
\[
 I_{\mathrm{lat}}(Y\times\PP^1)=2I_{\mathrm{lat}}(Y)>0.
\]
Every center in projective weak factorization of fourfolds has dimension
at most two, so Proposition~\ref{prop:atomic-lowdim} makes every correction
zero.  The intrinsic blowup formula therefore makes \(I_{\mathrm{lat}}\)
birationally invariant in dimension four.  The generic even connection of
\(\PP^4\) has only rank-one primary summands, so
\(I_{\mathrm{lat}}(\PP^4)=0\).  Hence \(Y\times\PP^1\) is irrational.
Rationality of \(Y\) would imply rationality of this product.
Finally, \(I_{\mathrm{exp}}\le I_{\mathrm{lat}}\) gives the last assertion.
\proves{cor:threefold-marker}%
\uses{prop:atomic-lowdim,prop:intrinsic-count-formulas}%
\imports{AKMW:weak-factorization}%
\end{proof}

